\documentclass[11pt]{article}
\usepackage{fontspec,amsmath,amssymb,amsthm,hyperref}
\usepackage[margin=1in]{geometry}
\setlength{\emergencystretch}{3em}
\title{Gibbs Dirichlet form, symmetry and dissipation}
\author{LangevinLean mathematical exposition}
\date{}
\begin{document}\maketitle
\noindent Source module: \texttt{UnadjustedLangevin/GeneratorDirichlet.lean}. Compile this proof document with XeLaTeX or LuaLaTeX. The accompanying HTML pairs the mathematical argument with the exact Lean statements and records the scope of the formalization.


For compactly supported \(C^2\) tests \(\varphi,\psi\), the generator obeys \(\int\varphi\mathcal L\psi\,d\pi=-\int\langle\nabla\varphi,\nabla\psi\rangle d\pi\). This provides the symmetric Dirichlet form at the test-function level.


\section{fderiv exp neg mul apply}


Differentiate the product \(e^{-U}\varphi\): \(D(e^{-U}\varphi)(x)[v]=e^{-U(x)}(D\varphi(x)[v]-\langle\nabla U(x),v\rangle\varphi(x))\). Lean combines the product rule with the already proved derivative of the Gibbs weight, then collects the common exponential factor algebraically.



\noindent Lean declarations: \texttt{fderiv\_exp\_neg\_mul\_apply}.

\section{integral weighted directional dirichlet}


Integrate by parts the derivative of \(e^{-U}\varphi\) against \(D\psi[v]\). The weighted product rule separates a \(D\varphi[v]D\psi[v]\) term from the drift term, giving \(\int e^{-U}\varphi L_v\psi=-\int e^{-U}D\varphi[v]D\psi[v]\). Compact support and \(C^2\) regularity justify every derivative and integral involved.



\noindent Lean declarations: \texttt{integral\_weighted\_directional\_dirichlet}.

\section{integral directionalGenerator dirichlet gibbs}


Convert both sides of the weighted directional identity to Gibbs integrals. The same real normalization factor multiplies both integrals, producing \(\int\varphi L_v\psi\,d\pi_U=-\int D\varphi[v]D\psi[v]d\pi_U\). This is a change of measure, with no new integration-by-parts step.



\noindent Lean declarations: \texttt{integral\_directionalGenerator\_dirichlet\_gibbs}.

\section{integrable gradient inner gibbsMeasure}


The gradients are continuous; the gradient of the compactly supported test \(\varphi\) vanishes outside its compact support. Thus \(x\mapsto\langle\nabla\varphi(x),\nabla\psi(x)\rangle\) is continuous with compact support. The general compact-test Gibbs integrability theorem applies even though \(\psi\) itself need not have compact support in this statement.



\noindent Lean declarations: \texttt{integrable\_gradient\_inner\_gibbsMeasure}.

\section{integrable test mul langevinGenerator gibbsMeasure}


The generator of the \(C^2\) test \(\psi\) is continuous under the Lipschitz-gradient assumption on \(U\). Multiplication by the continuous compactly supported \(\varphi\) makes the product compactly supported, hence Gibbs integrable. This is the exact product needed on the left side of the Dirichlet identity.



\noindent Lean declarations: \texttt{integrable\_test\_mul\_langevinGenerator\_gibbsMeasure}.

\section{integral langevinGenerator dirichlet gibbs}


Sum the directional Dirichlet identities over the orthonormal basis. Finite-sum integrability allows the integral to pass through the sum, and Parseval identifies \(\sum_iD\varphi[e_i]D\psi[e_i]=\langle\nabla\varphi,\nabla\psi\rangle\). Therefore \(\int\varphi L\psi\,d\pi_U=-\int\langle\nabla\varphi,\nabla\psi\rangle d\pi_U\).



\noindent Lean declarations: \texttt{integral\_langevinGenerator\_dirichlet\_gibbs}.

\section{integral langevinGenerator symmetric gibbs}


Apply the Dirichlet identity in both orders \((\varphi,\psi)\) and \((\psi,\varphi)\). Symmetry of the real inner product makes the resulting energy integrands equal pointwise. Hence \(\int\varphi L\psi\,d\pi_U=\int\psi L\varphi\,d\pi_U\).



\noindent Lean declarations: \texttt{integral\_langevinGenerator\_symmetric\_gibbs}.

\section{integral langevinGenerator self gibbs}


Specialize the bilinear identity to \(\psi=\varphi\), then rewrite \(\langle\nabla\varphi,\nabla\varphi\rangle=\|\nabla\varphi\|^2\). This identifies the quadratic generator form exactly: \(\int\varphi L\varphi\,d\pi_U=-\int\|\nabla\varphi\|^2d\pi_U\).



\noindent Lean declarations: \texttt{integral\_langevinGenerator\_self\_gibbs}.

\section{integral langevinGenerator dissipative gibbs}


The self-pairing identity expresses the generator quadratic form as minus an integral of a nonnegative square. The integral is nonnegative, so its negative is at most zero. The proof uses the actual integrable energy established earlier, giving dissipativity on compactly supported \(C^2\) tests.



\noindent Lean declarations: \texttt{integral\_langevinGenerator\_dissipative\_gibbs}.

\end{document}
